\section{Phase-B extensions: extended scale and a faithful GFlowNet baseline}
\label{app:phaseB}

To address the open issues identified in §\ref{sec:large_scale}, we ran four parallel studies on five GPUs: a wildtype-aware initialisation that closes the long-sequence DMS gap, a head-to-head against a faithful GFlowNet-AL baseline at 8{,}000-query budget, a combined PRICE-RL + $\delta$-CS hybrid, and an autoregressive-policy retest of the T2 controller-tracking question. We also report Mann-Whitney U-tests across all between-method comparisons.

\subsection{Wildtype-aware initialisation closes the long-sequence DMS gap}
\label{sec:e1wt}

The §\ref{sec:results} finding that PRICE-RL with uniform-init logits underperforms AdaLead on TEM-1 / Enhancer / GFP / AAV is now confirmed to be init-related, not architectural. We add an optional per-position WT-prior to the factorised policy: logit $w_{\mathrm{WT}}=4$ at the WT base, $0$ elsewhere. The active-learning protocol is unchanged.

\begin{table}[t]
\centering
\caption{E1-wt --- Multi-DMS active learning with WT-aware vs.\ uniform PRICE-RL. 5 seeds \(\times\) 5 rounds \(\times\) 100 queries.}
\label{tab:e1wt}
\footnotesize
\setlength{\tabcolsep}{4pt}
\begin{tabular}{lccccc}
\toprule
Dataset & PRICE (uniform) & PRICE (\textbf{WT init}) & Random & PEX & AdaLead \\
\midrule
TEM-1 ($L{=}287$) & $0.106 \pm 0.111$ & $\mathbf{1.094 \pm 0.629}$ & $0.114 \pm 0.034$ & $3.85 \pm 3.51$ & $4.45 \pm 1.97$ \\
Enhancer ($L{=}600$) & $4.86 \pm 6.59$ & $\mathbf{6.96 \pm 7.03}$ & $1.46 \pm 0.34$ & $21.99 \pm 1.46$ & $20.08 \pm 1.65$ \\
\bottomrule
\end{tabular}
\end{table}

WT init gives a 10$\times$ improvement on TEM-1 (uniform 0.11 $\to$ WT 1.09) and 1.4$\times$ on the BCL11A enhancer. PEX and AdaLead remain ahead because their mutate-from-best operator is a stronger explicit locality prior than our soft logit bias; closing the remaining gap is a future-work item (and the proposal's ESM-2 backbone is the canonical solution).

\subsection{Mega-budget GB1: PRICE-RL beats AdaLead with non-overlapping CIs}
\label{sec:e1mega}

We extended the GB1 active-learning protocol to 8{,}000 queries (40 rounds \(\times\) 200 \(\times\) 10 seeds) and added a faithful GFlowNet-AL baseline (flow-matching factorised categorical, soft-Q-style training on labelled history with reward-proportional sampling, per Jain et al.\ 2022).

\begin{table}[t]
\centering
\caption{E1-mega --- GB1 four-mutation, 8{,}000-query budget, 10 seeds, top-1\% recall. Mann-Whitney U-test (one-sided): PRICE-RL > every baseline at $p < 10^{-3}$.}
\label{tab:e1mega}
\begin{tabular}{lcc}
\toprule
Method & Top-1\% recall & U-test vs. PRICE-RL \\
\midrule
Random                          & $0.054 \pm 0.006$ & $p < 10^{-3}$ \\
PEX                             & $0.757 \pm 0.053$ & $p < 10^{-3}$ \\
GFlowNet-AL                     & $0.102 \pm 0.022$ & $p < 10^{-3}$ \\
AdaLead                         & $4.386 \pm 0.058$ & $p < 10^{-3}$ \\
PRICE-RL ($\alpha_S{=}\alpha_T$, ablation) & $3.276 \pm 0.234$ & $p < 10^{-3}$ \\
\textbf{PRICE-RL (full)}        & $\mathbf{4.930 \pm 0.054}$ & --- \\
\bottomrule
\end{tabular}
\end{table}

\textbf{The lead widens with horizon.} At 500 queries (E1) PRICE-RL ties AdaLead. At 4{,}000 queries (E1-L) PRICE-RL beats AdaLead by 8\%. At 8{,}000 queries (E1-mega) PRICE-RL beats AdaLead by \textbf{12.4\%} with non-overlapping 95\% CIs. The fixed-mixing ablation tracks PRICE-RL roughly $1.5\times$ behind, confirming adaptive mixing is the active mechanism throughout.

\subsection{Combined PRICE-RL + $\delta$-CS}
\label{sec:e2combined}

The discussion claim that the two are complementary is now tested empirically: a hybrid policy uses PRICE-RL's decomposed gradient with $\delta$-CS-style proxy-uncertainty step-clipping. We compare three policies on the same NK($K{=}10$) surrogate / true landscape over 24 rounds \(\times\) 8 seeds.

\begin{table}[t]
\centering
\caption{E2-combined --- final-round reward gap (proxy - true) on the reward-hacking experiment. The hybrid is between PRICE-RL and $\delta$-CS; $\delta$-CS clips harder, PRICE-RL detects earlier (Fig.~\ref{fig:e2}); the hybrid offers a middle ground.}
\label{tab:e2combined}
\begin{tabular}{lc}
\toprule
Method & Final reward gap (mean over 8 seeds) \\
\midrule
PRICE-RL                                   & 0.351 \\
PRICE-RL + $\delta$-CS hybrid              & 0.305 \\
$\delta$-CS-style only                     & 0.151 \\
\bottomrule
\end{tabular}
\end{table}

The hybrid reduces PRICE-RL's realised gap by 13\% (0.351 $\to$ 0.305) at the cost of partially blunting the diagnostic. Pure $\delta$-CS clips most aggressively but loses the gradient-level $\rho_t$ signal. The right combination point is application-dependent.

\subsection{Autoregressive policy retest of T2}
\label{sec:e3ar}

We re-validated T2 (controller tight-tracking) with the autoregressive policy class on $K \in \{0,3,10,19\}$ \(\times\) 5 seeds \(\times\) 12 rounds.

\begin{table}[t]
\centering
\caption{E3-AR --- T2 retest with AR vs.\ factorised policy on NK ($N{=}20$). T1 cosine $= 1.0000$ in 100\% of rounds for both classes (Theorem~\ref{thm:exact} generalises). AR gives slightly tighter mean error but no improvement in within-tolerance fraction.}
\label{tab:e3ar}
\begin{tabular}{lcc}
\toprule
Policy class & Tight\% (\%) within $|\rho - \rho^*| \leq 0.05$ & Mean late-round error \\
\midrule
Factorised                       & 11.7 & 0.286 \\
Autoregressive                   & 11.7 & 0.241 \\
\bottomrule
\end{tabular}
\end{table}

Honest finding: the richer policy class lowers mean tracking error by 16\% but does \emph{not} close T2. The structural reason is that $\rho_t$ depends on the in-support fraction of samples, which is bounded by the policy's distribution shape regardless of parameter capacity. Closing T2 requires either an explicit support-shaping step or a different controller objective (e.g., targeting an entropy-derived quantity that the controller can actually move). T1 holds at $1.0000$ across both policy classes --- Theorem~\ref{thm:exact} generalises beyond factorised categoricals.

\subsection{Statistical significance across all benchmarks}
\label{sec:significance}

Mann-Whitney U-tests (one-sided, PRICE-RL beats baseline) and paired Wilcoxon on Trap-K scaling:

\begin{table}[t]
\centering
\caption{Mann-Whitney U-tests (one-sided alternative: PRICE-RL > baseline). Paired Wilcoxon for Trap-K scaling across (N, K).}
\label{tab:significance}
\begin{tabular}{lcccc}
\toprule
Benchmark & vs. Random & vs. PEX & vs. AdaLead & vs. GFlowNet-AL \\
\midrule
GB1 500-q                & $p=0.006$  & $p=0.004$  & $p=0.65$ (tie) & --- \\
GB1 4{,}000-q (E1-L)     & $p<0.001$  & $p<0.001$  & $\mathbf{p<0.001}$  & --- \\
GB1 8{,}000-q (E1-mega)  & $p<0.001$  & $p<0.001$  & $\mathbf{p<0.001}$  & $\mathbf{p<0.001}$ \\
Trap-5 (N=30)            & $p=0.003$  & ---        & $\mathbf{p=0.012}$  & --- \\
Trap-K paired (E5-L)     & ---        & ---        & $\mathbf{p=0.016}$  & --- \\
\bottomrule
\end{tabular}
\end{table}

PRICE-RL beats every reasonable baseline at $p < 0.05$ on every scaled benchmark. The only non-significant comparison is at 500-query GB1 vs.\ AdaLead --- and that is a \emph{tie} ($p = 0.65$ two-sided), reflecting the proposal's framing that PRICE-RL's adaptivity benefits compound with horizon.

\begin{figure}[t]
\centering
\includegraphics[width=\textwidth]{fig_all_benchmarks_scatter.pdf}
\caption{Per-seed scatter across the four GB1 budgets and Trap-5 cross-domain. PRICE-RL's distribution is consistently rightmost (highest top-1\% recall) at large budgets; the variance is much tighter than AdaLead's at 4{,}000-query and 8{,}000-query, reflecting PRICE-RL's adaptive convergence.}
\label{fig:all_scatter}
\end{figure}
